\documentclass[12pt]{article}
\usepackage[margin=1in]{geometry}
\usepackage{enumitem}
\usepackage{titlesec}
\usepackage{parskip}
\usepackage{graphicx}
\usepackage{xcolor}
\usepackage{hyperref}
\usepackage{fontenc}

\title{\textbf{Gandhi and Przeworski 2007 RG Notes}\\
\large Gandhi, Jennifer, and Adam Przeworski. ``Authoritarian Institutions and the Survival of Autocrats'' [2007].\\
\textit{Comparative Political Studies} 40(11): 1279-1301.}
\author{}
\date{}

\begin{document}

\maketitle

Provides a basis for the inclusion of non-unanimous legislatures even within authoritarian regimes -- such structures allow opposition forces to be absorbed and lend some tempered amount of credible power to them. Including opposition forces allows the leader to lengthen their tenure

Relate this back to mandates to rule (Powell)

Specifically argues that the outcome is length -- but, to some extent, perhaps there is a degree of latitude achievable as well: if the autocrat can survive even under some degree of opposition, then they surely must have been granted a mandate to rule\ldots{}

``In turn, although not necessarily universalistic, policy concessions must be formalized as legal norms. Hence, we believe that although spoils can be distributed directly out of the autocrat's pocket, working out policy concessions requires an institutional setting: some forum to which access can be controlled, where demands can be revealed without appearing as acts of resistance, where compromises can be hammered out without undue public scrutiny, and where the resulting agreements can be dressed in a legalistic form and publicized as such.'' (Gandhi and Przeworski, 2007, p. 1282)

Extremely similar to North \& Weingast, and A\&R -- extend political enfranchisement as a credible commitment and in doing so sustain for longer an equilibrium whereby power is maintained

Important consideration: when a nation is endowed with resources, we can expect fewer legislative parties -- this is because the government does not actually need cooptation of the citizens in order to extract wealth, and can simply manage that on its own

\end{document}
